\section{Code availability and repository structure}
\label{app:code-availability}

The implementation used in this dissertation is available at
\url{https://github.com/Linlin021222/mental-health-support-crisis-chats}.
The repository contains the source code for preprocessing, model training,
user-disjoint validation, prediction, calibration and the Kaggle Qwen
experiments. The competition data are not included because access to the UMD
Reddit data is restricted. Model checkpoints, cached probabilities and
leaderboard submissions are also excluded.

The main local entry point is \texttt{main.py}. Project settings and label
definitions are stored in \texttt{configs/config.py}. Data processing is in
\texttt{preprocess/} and \texttt{datasets/}. Neural architectures are in
\texttt{models/}, while training and validation procedures are in
\texttt{trainer/}. Prediction and ensemble logic are stored in
\texttt{inference/}. The classic BERT comparison is implemented in
\texttt{baselines/bert\_multitask\_baseline.py}.

The Task 1 Qwen3-8B QLoRA experiment is stored in
\texttt{kaggle\_llm\_package/kaggle\_llm\_experiment.py}. The final Task 2
pairwise model is stored in
\texttt{kaggle\_factor\_reranker\_package/kaggle\_factor\_reranker\_v3.py},
with its factor definitions in \texttt{factor\_taxonomy\_v3.py}. Each Kaggle
folder contains a separate README describing fold evaluation and full-data
training.

To reproduce the local pipeline, an authorised user should place
\texttt{train.xlsx} and \texttt{leaderboard.xlsx} in the \texttt{data/}
directory and install the packages listed in \texttt{requirements.txt}. The
command \texttt{python main.py --mode private-group-cv} runs the grouped local
evaluation. Full-data training and prediction are started with
\texttt{python main.py --mode train-full} and
\texttt{python main.py --mode predict}. The exact final ensemble also requires
the Qwen adapters and probability files generated on Kaggle. These files are
not distributed in the public repository.
